\begin{frame}{Pensiones}
\begin{itemize}
    \item El gasto pensional incluye a Colpensiones y a los regímenes especiales de las Fuerzas Militares, los docentes y otros empleados públicos.
    \item Entre 2024 y 2035, el gasto en pensiones podría aumentar cerca de 1 punto del PIB.
    \item La cobertura sigue siendo baja: solo uno de cada cuatro adultos mayores recibe pensión.
\end{itemize}
\end{frame}

\begin{frame}{Pensiones}
\begin{enumerate}
    \item Ajustes paramétricos:
    \begin{itemize}
        \item edad de pensión;
        \item tasa de reemplazo;
        \item umbral de cotización.
    \end{itemize}
    \item Eliminar la exención de impuestos a las pensiones.
    \item Aumentar la cotización en salud de los pensionados.
    \item Incrementar los aportes de los pensionados al Fondo de Solidaridad Pensional.
\end{enumerate}
\end{frame}

\begin{frame}{Pensiones}
\begin{itemize}
    \item La Ley 2381 de 2024 dirige más cotizaciones al componente público y crea un fondo de ahorro administrado por el Banco de la República.
    \item El uso de los recursos del fondo ante presiones fiscales adelantaría su agotamiento.
    \item Los ajustes al umbral de ingresos, la tasa de reemplazo y la edad de jubilación deben evaluarse por cohortes.
    \item Una reforma temprana reduciría el pasivo implícito y enviaría una señal de compromiso fiscal.
\end{itemize}
\end{frame}